\section{Forest subgroups and their incidence geometry}
Let $T$ be a finite tree on a set $\Omega$ of $N\geq2$ vertices, with edge
set $E$. For $F\subseteq E$, let $K_F$ be the subgroup of $\Sym(\Omega)$
generated by the endpoint transpositions of the edges in $F$. Set
$H_e=K_{E\setminus\{e\}}$. The elements of type $e$ in $\Gamma(T)$ are
left cosets $gH_e$. Two such elements are incident when their cosets
intersect; incidence is reflexive, and distinct elements of one type are
not incident. A flag is a set of pairwise incident elements. The residue
of a flag consists of its incident elements of the remaining types.

\begin{lemma}\label{forest}
If $C$ ranges over the components of $(\Omega,F)$, then
\[
 K_F=\prod_C\Sym(C),\qquad K_F\cap K_D=K_{F\cap D}.
\]
The coset system $\Gamma(T)$ is a thin, residually connected geometry on
which $\Sym(\Omega)$ is flag-transitive. The residue at the standard flag
of types $E\setminus F$ is the join of the geometries of the nontrivial
components of $(\Omega,F)$.
\end{lemma}
\begin{proof}
Transpositions along a connected graph generate its full symmetric group:
conjugating successive edge transpositions along a path produces the
transposition of its endpoints. Distinct components have disjoint
supports. A permutation in both $K_F$ and $K_D$ preserves every block of
the common refinement of their component partitions. Two vertices are
in such a block exactly when their unique tree path belongs to both
forests, hence to $F\cap D$. This proves the two group formulas.

Root $T$, and let $A_e$ be the component of $T-e$ away from the root.
The coset $gH_e$ identifies with the subset $gA_e$ of $\Omega$.
For distinct types, incidence is precisely
\[
 |X_e\cap X_f|=|A_e\cap A_f|.
\]
Indeed these intersection sizes and the two fixed subset sizes determine
all four atoms of the pair; equal atom sizes give a permutation carrying
one pair to the other. The $A_e$ form a laminar family. Thus pairwise
incidence forces exactly their nesting and disjointness relations on the
$X_e$. Adjoin $\Omega$ above the selected sets and subtract each
container's maximal selected proper subcontainers. The resulting atoms
partition $\Omega$ and have the same sizes on the standard and target
sides. Bijections of corresponding atoms combine to one permutation
realizing the entire flag. Every flag therefore extends to a chamber,
and flags of each fixed type form one orbit.

The stabilizer of the standard flag of types $E\setminus F$ is $K_F$.
Its maximal parabolics for the remaining types are $K_{F\setminus\{e\}}$.
Their direct product descriptions identify the residue with the claimed
join; isolated vertices contribute trivial factors. The chamber
stabilizer is $K_\varnothing=1$, and each panel stabilizer has order two,
so the geometry is thin. In any residue of rank at least two, the
maximal parabolics collectively contain every remaining edge
transposition, and hence generate the residue group. Their coset
incidence graph is connected: products of subgroup elements give paths
between standard and arbitrary cosets. This proves residual connectedness.
\end{proof}

For an element type $i$, write $\sigma_i(F)$ for the set of elements of
type $i$ incident with every member of $F$. The \emph{intersection
property} requires that for every element $x$ and flag $F$, a nonempty
$\sigma_i(x)\cap\sigma_i(F)$ be $\sigma_i(F')$ for a flag $F'$ incident
with $x$ and $F$. This is required in the geometry and in every residue.
It is distinct from the subgroup intersection identity of Lemma~\ref{forest}.

\section{Explicit witnesses for Boolean shadows}
Let $\mathcal B(\Omega)$ have the nonempty proper subsets of $\Omega$ as
elements, cardinality as type, and comparability by inclusion as incidence.

\begin{lemma}\label{boolean}
Fix $0<k<N$. For any family $\mathcal U$ of elements of
$\mathcal B(\Omega)$, a nonempty common $k$-shadow has a witness flag of
at most two elements, incident with every member of $\mathcal U$.
\end{lemma}
\begin{proof}
If $\mathcal U$ contains a $k$-subset $X$, every common shadow element
must equal $X$. Nonemptiness says that $X$ is incident with the whole
family, so $\{X\}$ is the required flag. Otherwise put
\[
 L=\bigcup_{\substack{X\in\mathcal U\\|X|<k}}X,
 \qquad V=\bigcap_{\substack{X\in\mathcal U\\|X|>k}}X,
\]
with empty union $\varnothing$ and empty intersection $\Omega$.
The common shadow is exactly
\begin{equation}\label{interval}
 \{Y\subseteq\Omega: |Y|=k,\ L\subseteq Y\subseteq V\}.
\end{equation}
Nonemptiness gives $L\subseteq V$ and $|L|\leq k\leq|V|$.
If either cardinality bound is equality, the shadow has a unique element,
which again supplies a singleton flag incident with the whole family.
In the remaining case use $\{L,V\}$, omitting $\varnothing$ and
$\Omega$. Its members are proper subsets of different sizes and form
a flag. Every smaller query member is contained in $L$, and every
larger one contains $V$, so both witness elements are comparable with
every query member. Its shadow is precisely~\eqref{interval}, including
when one or both boundary sets were omitted.
\end{proof}

\begin{lemma}\label{boolean-residues}
The geometry $\mathcal B(\Omega)$ satisfies the intersection property
in all residues. Joins of geometries with that property also satisfy it.
\end{lemma}
\begin{proof}
Apply Lemma~\ref{boolean} to $\mathcal U=\{x\}\cup F$ for the property
in $\mathcal B(\Omega)$. A fixed chain, completed by the formal endpoints
$\varnothing$ and $\Omega$, divides its residue into the proper intervals
between successive members. An interval $(A,B)$ is the Boolean geometry
on $B\setminus A$, via $X\mapsto X\setminus A$; the types are shifted
by $|A|$. All elements in different intervals are incident, so the residue
is their join. Empty intervals add no types. A rank-one factor has only
empty, full or singleton shadows; the empty flag witnesses a full
shadow, and its unique point witnesses a singleton.

For completeness, in a join fix a type $i$ in one factor. Only query
elements of that factor constrain its shadow. If $x$ belongs to it, use
the factor's intersection property with the part of $F$ in that factor.
If $x$ belongs to another factor, the part of $F$ in the target factor
is itself a witness, including when it is empty. The witness is
automatically incident with all query elements in other factors.
Residues of a join are joins of factor residues, proving the assertion
also at every residue.
\end{proof}

\begin{theorem}\label{path}
The geometry $\Gamma(T)$ satisfies the intersection property if and only
if $T$ is a path.
\end{theorem}
\begin{proof}
For a path, root at one endpoint. Its cuts $A_e$ form a chain with sizes
$1,\ldots,N-1$. The subset description in Lemma~\ref{forest} identifies
its incidence with comparability, so $\Gamma(T)=\mathcal B(\Omega)$.
Lemma~\ref{boolean-residues} proves the positive assertion.

A tree which is not a path has three edges incident at a common vertex.
Keep precisely those edges in Lemma~\ref{forest}. The resulting residue
is the geometry of the claw, with three types $a,b,c$, each a copy of
$\{1,2,3,4\}$. Elements of different types are incident exactly when
their labels differ. Take $x=(b,1)$ and $F=\{(b,2)\}$, with target
type $a$. Their common shadow is $\{(a,3),(a,4)\}$. A witness cannot
contain an $a$-element, which would give at most a singleton shadow,
or a $b$-element, which cannot be incident with both distinct
$b$-elements. It can contain at most one $c$-element. The empty flag
has four shadow elements, whereas any $c$-element has three. Neither
has the required shadow. Thus this residue fails the property.
\end{proof}

